\documentclass[11pt]{article}
\usepackage[provide=*,spanish]{babel}
\usepackage{fontspec}
\defaultfontfeatures{Ligatures=TeX}
\setmainfont{Latin Modern Roman}
\setsansfont{Latin Modern Sans}
\setmonofont{Latin Modern Mono}
\usepackage[a4paper,margin=2.25cm,headheight=15pt]{geometry}
\usepackage{amsmath,amssymb}
\usepackage{booktabs,longtable,tabularx,array}
\usepackage{graphicx,float}
\usepackage{xcolor}
\usepackage{hyperref}

\definecolor{uvgverde}{HTML}{176B52}
\definecolor{uvgoscuro}{HTML}{103F33}
\definecolor{grisclaro}{HTML}{F2F5F4}
\definecolor{azulcli}{HTML}{17324D}
\hypersetup{colorlinks=true,linkcolor=uvgoscuro,urlcolor=uvgverde,pdfauthor={Grupo 1 - CC3067},pdftitle={Laboratorio 3 - Algoritmos de Enrutamiento}}

\pagestyle{plain}
\renewcommand{\arraystretch}{1.18}

\newcolumntype{Y}{>{\raggedright\arraybackslash}X}
\newcolumntype{C}[1]{>{\centering\arraybackslash}p{#1}}
\newenvironment{calculo}[1][]{\par\medskip\noindent\color{uvgoscuro}\hrule\smallskip\textbf{Desarrollo del calculo}\par\color{black}}{\par\smallskip\color{uvgoscuro}\hrule\color{black}\medskip}
\newenvironment{decision}[1][]{\par\medskip\noindent\color{uvgoscuro}\hrule\smallskip\textbf{Decision de diseno}\par\color{black}}{\par\smallskip\color{uvgoscuro}\hrule\color{black}\medskip}
\newenvironment{evidencia}[1][]{\par\medskip\noindent\color{azulcli}\hrule\smallskip\textbf{Evidencia real de Packet Tracer}\par\color{black}}{\par\smallskip\color{azulcli}\hrule\color{black}\medskip}

\newcommand{\universidad}{Universidad del Valle de Guatemala}
\newcommand{\curso}{CC3067 Redes}
\newcommand{\seccion}{Seccion 10}
\newcommand{\grupo}{Grupo 1}
\newcommand{\integrantes}{%
\begin{tabular}{ll}
Pablo Daniel Barillas Moreno & 22193\\
Andres Rafael Chivalan & 21534
\end{tabular}}

\newcommand{\portadalab}[2]{%
\begin{titlepage}
\centering
{\Large\bfseries \universidad\par}
\vspace{1.2cm}
{\large \curso\ -- \seccion\par}
\vspace{2.2cm}
{\Huge\bfseries Laboratorio 3\par}
\vspace{0.35cm}
{\LARGE Algoritmos de Enrutamiento\par}
\vspace{0.8cm}
{\Large\color{uvgverde}#1\par}
\vfill
{\large\bfseries \grupo\par}
\vspace{0.45cm}
\integrantes
\vfill
{\large #2\par}
\end{titlepage}}



\begin{document}
\portadalab{PDF 1: desarrollo manual de VLSM y conceptos}{5 de octubre de 2026}
\tableofcontents
\newpage

\section{Alcance y criterio usado}

Este documento funciona como base para transcribir a mano el procedimiento. No se limita a presentar resultados: en cada subred se muestra la cantidad de bits de host, el prefijo, la mascara, el tamano del bloque, las direcciones utilizables, el broadcast y el desperdicio.

El documento principal del laboratorio es la fuente de requisitos. Los documentos de direccionamiento, enrutamiento dinamico y configuracion OSPF se utilizaron como apoyo teorico. Por indicacion adicional del estudiante, los pasos se presentan con sustituciones numericas y comprobaciones explicitas.

\begin{decision}[title={Aclaracion sobre R-EIGRP-1}]
La tabla general del enunciado exige ``3 VLANs'', mientras que la tabla detallada solo asigna proposito y cantidad de hosts a VLAN 10 (60 usuarios) y VLAN 20 (30 dispositivos de voz). Para satisfacer la condicion mas restrictiva se implemento una tercera VLAN 30 de administracion/servidores con 10 hosts, usando el siguiente bloque VLSM disponible. Esta decision no modifica los segmentos definidos: agrega una red pequena y coherente con el rol de un sitio de acceso.
\end{decision}

\section{Bloques base y estrategia VLSM}

\begin{center}
\begin{tabularx}{\textwidth}{l l Y}
\toprule
Dominio & Bloque base & Motivo\\
\midrule
LAN/VLAN OSPF & \texttt{192.168.0.0/20} & Direccion privada de clase C solicitada, con 4096 direcciones totales.\\
LAN/VLAN EIGRP & \texttt{172.16.0.0/20} & Direccion privada de clase B solicitada, con 4096 direcciones totales.\\
Enlaces seriales & \texttt{10.0.0.0/24} & Direccion privada de clase A solicitada; se divide en seis subredes /30.\\
\bottomrule
\end{tabularx}
\end{center}

Para una red con $H$ hosts se busca el menor numero entero $h$ que cumpla
\[
2^h-2\ge H.
\]
Se restan dos direcciones porque una identifica la red y otra es el broadcast. El prefijo se calcula como
\[
p=32-h.
\]
El tamano o salto del bloque es $2^h$ direcciones. Las solicitudes se asignan de mayor a menor para evitar fragmentacion.

\section{Desarrollo VLSM del dominio OSPF}

\subsection{LAN de R-OSPF-2: 200 usuarios}
\begin{calculo}
\[
2^7-2=126<200,\qquad 2^8-2=254\ge 200.
\]
Se necesitan $h=8$ bits de host. Por tanto,
\[
p=32-8=24.
\]
La mascara /24 es \texttt{255.255.255.0}. El bloque contiene $2^8=256$ direcciones y avanza de 256 en el ultimo octeto.

Primera subred disponible: \texttt{192.168.0.0/24}.
\begin{align*}
\text{Red} &= \texttt{192.168.0.0},\\
\text{Primer host} &= \texttt{192.168.0.1},\\
\text{Ultimo host} &= \texttt{192.168.0.254},\\
\text{Broadcast} &= \texttt{192.168.0.255},\\
\text{Desperdicio} &= 254-200=54\text{ hosts}.
\end{align*}
\end{calculo}

\subsection{VLAN 10 de R-OSPF-1: 50 usuarios}
\begin{calculo}
\[
2^5-2=30<50,\qquad 2^6-2=62\ge 50.
\]
Se necesitan $h=6$ bits y $p=32-6=26$. La mascara es
\[
11111111.11111111.11111111.11000000=\texttt{255.255.255.192}.
\]
El salto es $256-192=64$. Despues del /24 anterior, la siguiente red libre es \texttt{192.168.1.0/26}.
\begin{align*}
\text{Red} &= \texttt{192.168.1.0}, & \text{Primer host} &= \texttt{192.168.1.1},\\
\text{Ultimo host} &= \texttt{192.168.1.62}, & \text{Broadcast} &= \texttt{192.168.1.63},\\
\text{Desperdicio} &= 62-50=12\text{ hosts}.&&
\end{align*}
\end{calculo}

\subsection{VLAN 30 de R-OSPF-1: 10 hosts administrativos/servidores}
\begin{calculo}
\[
2^3-2=6<10,\qquad 2^4-2=14\ge 10.
\]
Se necesitan $h=4$ bits, entonces $p=28$. La mascara es \texttt{255.255.255.240} y el salto es $256-240=16$. La red comienza inmediatamente despues del broadcast anterior:
\begin{align*}
\text{Red} &= \texttt{192.168.1.64/28},\\
\text{Primer host} &= \texttt{192.168.1.65},\\
\text{Ultimo host} &= \texttt{192.168.1.78},\\
\text{Broadcast} &= \texttt{192.168.1.79},\\
\text{Desperdicio} &= 14-10=4\text{ hosts}.
\end{align*}
\end{calculo}

\subsection{Resumen del VLSM OSPF}
\begin{center}
\small
\begin{tabularx}{\textwidth}{Y C{1.25cm} C{1.35cm} l l l C{1.25cm}}
\toprule
Segmento & Hosts & Prefijo & Red & Primer host & Ultimo host / broadcast & Libres\\
\midrule
R-OSPF-2 usuarios & 200 & /24 & 192.168.0.0 & 192.168.0.1 & .254 / .255 & 54\\
R-OSPF-1 VLAN 10 & 50 & /26 & 192.168.1.0 & 192.168.1.1 & .62 / .63 & 12\\
R-OSPF-1 VLAN 30 & 10 & /28 & 192.168.1.64 & 192.168.1.65 & .78 / .79 & 4\\
\bottomrule
\end{tabularx}
\end{center}

\section{Desarrollo VLSM del dominio EIGRP}

\subsection{LAN de R-EIGRP-2: 100 usuarios}
\begin{calculo}
\[
2^6-2=62<100,\qquad 2^7-2=126\ge100.
\]
$h=7$, $p=25$, mascara \texttt{255.255.255.128} y salto 128.
\begin{align*}
\text{Red} &= \texttt{172.16.0.0/25}, & \text{Primer host} &= \texttt{172.16.0.1},\\
\text{Ultimo host} &= \texttt{172.16.0.126}, & \text{Broadcast} &= \texttt{172.16.0.127},\\
\text{Desperdicio} &=126-100=26.&&
\end{align*}
\end{calculo}

\subsection{VLAN 10 de R-EIGRP-1: 60 usuarios}
\begin{calculo}
\[
2^5-2=30<60,\qquad 2^6-2=62\ge60.
\]
$h=6$, $p=26$, mascara \texttt{255.255.255.192} y salto 64. La siguiente frontera valida despues de .127 es .128.
\begin{align*}
\text{Red} &= \texttt{172.16.0.128/26}, & \text{Primer host} &= \texttt{172.16.0.129},\\
\text{Ultimo host} &= \texttt{172.16.0.190}, & \text{Broadcast} &= \texttt{172.16.0.191},\\
\text{Desperdicio} &=62-60=2.&&
\end{align*}
\end{calculo}

\subsection{VLAN 20 de voz: 30 dispositivos}
\begin{calculo}
\[
2^5-2=30\ge30.
\]
$h=5$, $p=27$, mascara \texttt{255.255.255.224} y salto $256-224=32$.
\begin{align*}
\text{Red} &= \texttt{172.16.0.192/27}, & \text{Primer host} &= \texttt{172.16.0.193},\\
\text{Ultimo host} &= \texttt{172.16.0.222}, & \text{Broadcast} &= \texttt{172.16.0.223},\\
\text{Desperdicio} &=30-30=0.&&
\end{align*}
\end{calculo}

\subsection{VLAN 30 de administracion/servidores: 10 hosts}
\begin{calculo}
\[
2^3-2=6<10,\qquad 2^4-2=14\ge10.
\]
Se requieren $h=4$ bits de host, por lo que $p=32-4=28$. La mascara es \texttt{255.255.255.240}; el salto es $256-240=16$. El primer limite disponible despues de \texttt{172.16.0.223} es .224:
\begin{align*}
\text{Red} &= \texttt{172.16.0.224/28}, & \text{Primer host} &= \texttt{172.16.0.225},\\
\text{Ultimo host} &= \texttt{172.16.0.238}, & \text{Broadcast} &= \texttt{172.16.0.239},\\
\text{Desperdicio} &=14-10=4.&&
\end{align*}
La IP \texttt{172.16.0.225} se asigno al gateway y \texttt{172.16.0.226} al equipo de prueba.
\end{calculo}

\subsection{Resumen del VLSM EIGRP}
\begin{center}
\small
\begin{tabularx}{\textwidth}{Y C{1.25cm} C{1.35cm} l l l C{1.25cm}}
\toprule
Segmento & Hosts & Prefijo & Red & Primer host & Ultimo host / broadcast & Libres\\
\midrule
R-EIGRP-2 usuarios & 100 & /25 & 172.16.0.0 & 172.16.0.1 & .126 / .127 & 26\\
R-EIGRP-1 VLAN 10 & 60 & /26 & 172.16.0.128 & 172.16.0.129 & .190 / .191 & 2\\
R-EIGRP-1 VLAN 20 & 30 & /27 & 172.16.0.192 & 172.16.0.193 & .222 / .223 & 0\\
R-EIGRP-1 VLAN 30 & 10 & /28 & 172.16.0.224 & 172.16.0.225 & .238 / .239 & 4\\
\bottomrule
\end{tabularx}
\end{center}

El intervalo \texttt{172.16.0.240--172.16.0.255} permanece disponible para crecimiento. Las cuatro redes quedan contenidas dentro del resumen \texttt{172.16.0.0/24}.

\section{Enlaces punto a punto}

Cada enlace requiere exactamente dos IP utilizables:
\[
2^1-2=0<2,\qquad 2^2-2=2\ge2.
\]
Por tanto, $h=2$, $p=30$, mascara \texttt{255.255.255.252}, bloque de cuatro direcciones y salto 4.

\begin{center}
\small
\begin{longtable}{p{4.0cm}p{2.7cm}p{2.7cm}p{2.7cm}}
\toprule
Enlace & Red /30 & Extremo A & Extremo B\\
\midrule
R-OSPF-1 -- R-OSPF-2 & 10.0.0.0 & 10.0.0.1 & 10.0.0.2\\
R-OSPF-1 -- R-CENTRAL & 10.0.0.4 & 10.0.0.5 & 10.0.0.6\\
R-OSPF-2 -- R-CENTRAL & 10.0.0.8 & 10.0.0.9 & 10.0.0.10\\
R-EIGRP-1 -- R-EIGRP-2 & 10.0.0.12 & 10.0.0.13 & 10.0.0.14\\
R-EIGRP-1 -- R-CENTRAL & 10.0.0.16 & 10.0.0.17 & 10.0.0.18\\
R-EIGRP-2 -- R-CENTRAL & 10.0.0.20 & 10.0.0.21 & 10.0.0.22\\
\bottomrule
\end{longtable}
\end{center}

Ejemplo para \texttt{10.0.0.16/30}: red .16, hosts .17 y .18, broadcast .19. El siguiente bloque comienza en .20.

\section{Tabla completa de direccionamiento}
\small
\begin{longtable}{p{2.6cm}p{3.1cm}p{2.8cm}p{2.8cm}p{2.7cm}}
\toprule
Dispositivo & Interfaz & Direccion & Mascara & Proposito\\
\midrule
\endhead
R-OSPF-1 & G0/0.10 & 192.168.1.1 & 255.255.255.192 & Gateway VLAN 10\\
R-OSPF-1 & G0/0.30 & 192.168.1.65 & 255.255.255.240 & Gateway VLAN 30\\
R-OSPF-1 & S0/0/0 & 10.0.0.1 & 255.255.255.252 & Hacia R-OSPF-2\\
R-OSPF-1 & S0/0/1 & 10.0.0.5 & 255.255.255.252 & Hacia central\\
R-OSPF-2 & G0/0 & 192.168.0.1 & 255.255.255.0 & Gateway usuarios\\
R-OSPF-2 & S0/0/0 & 10.0.0.2 & 255.255.255.252 & Hacia R-OSPF-1\\
R-OSPF-2 & S0/0/1 & 10.0.0.9 & 255.255.255.252 & Hacia central\\
R-CENTRAL & S0/0/0 & 10.0.0.6 & 255.255.255.252 & Hacia R-OSPF-1\\
R-CENTRAL & S0/0/1 & 10.0.0.10 & 255.255.255.252 & Hacia R-OSPF-2\\
R-CENTRAL & S0/1/0 & 10.0.0.18 & 255.255.255.252 & Hacia R-EIGRP-1\\
R-CENTRAL & S0/1/1 & 10.0.0.22 & 255.255.255.252 & Hacia R-EIGRP-2\\
R-EIGRP-1 & G0/0.10 & 172.16.0.129 & 255.255.255.192 & Gateway VLAN 10\\
R-EIGRP-1 & G0/0.20 & 172.16.0.193 & 255.255.255.224 & Gateway VLAN 20\\
R-EIGRP-1 & G0/0.30 & 172.16.0.225 & 255.255.255.240 & Gateway VLAN 30\\
R-EIGRP-1 & S0/0/0 & 10.0.0.13 & 255.255.255.252 & Hacia R-EIGRP-2\\
R-EIGRP-1 & S0/0/1 & 10.0.0.17 & 255.255.255.252 & Hacia central\\
R-EIGRP-2 & G0/0 & 172.16.0.1 & 255.255.255.128 & Gateway usuarios\\
R-EIGRP-2 & S0/0/0 & 10.0.0.14 & 255.255.255.252 & Hacia R-EIGRP-1\\
R-EIGRP-2 & S0/0/1 & 10.0.0.21 & 255.255.255.252 & Hacia central\\
\midrule
PC-O1-USUARIOS & FastEthernet0 & 192.168.1.2 & /26; GW 192.168.1.1 & Prueba VLAN 10\\
PC-O1-ADMIN & FastEthernet0 & 192.168.1.66 & /28; GW 192.168.1.65 & Prueba VLAN 30\\
PC-O2-USUARIOS & FastEthernet0 & 192.168.0.2 & /24; GW 192.168.0.1 & Prueba LAN\\
PC-E1-USUARIOS & FastEthernet0 & 172.16.0.130 & /26; GW 172.16.0.129 & Prueba VLAN 10\\
PC-E1-VOZ & FastEthernet0 & 172.16.0.194 & /27; GW 172.16.0.193 & Prueba VLAN 20\\
PC-E1-ADMIN & FastEthernet0 & 172.16.0.226 & /28; GW 172.16.0.225 & Prueba VLAN 30\\
PC-E2-USUARIOS & FastEthernet0 & 172.16.0.2 & /25; GW 172.16.0.1 & Prueba LAN\\
\bottomrule
\end{longtable}
\normalsize

\section{Calculo de las rutas resumidas}

\subsection{Resumen del dominio OSPF}
Los tres prefijos LAN ocupan desde \texttt{192.168.0.0} hasta \texttt{192.168.1.79}. Los terceros octetos relevantes son
\[
0=\texttt{00000000}_2,\qquad 1=\texttt{00000001}_2.
\]
Comparten los primeros siete bits del tercer octeto. Sumados a los 16 bits completos de \texttt{192.168}, existen $16+7=23$ bits comunes:
\[
\boxed{\texttt{192.168.0.0/23}}\qquad
\text{mascara }\texttt{255.255.254.0}.
\]
El rango del resumen es \texttt{192.168.0.0--192.168.1.255}; contiene las tres subredes.

\subsection{Resumen del dominio EIGRP}
Las subredes \texttt{172.16.0.0/25}, \texttt{172.16.0.128/26}, \texttt{172.16.0.192/27} y \texttt{172.16.0.224/28} comparten exactamente los primeros 24 bits:
\[
\boxed{\texttt{172.16.0.0/24}}\qquad
\text{mascara }\texttt{255.255.255.0}.
\]

\section{VLAN, enlace troncal y router-on-a-stick}

Una VLAN crea un dominio de broadcast logico. Los puertos de las computadoras son de acceso y transportan una sola VLAN. El enlace switch--router es troncal 802.1Q y transporta varias VLAN etiquetadas. El router crea una subinterfaz por VLAN; cada subinterfaz es el gateway de su red.

Ejemplo de R-OSPF-1:
\begin{verbatim}
interface GigabitEthernet0/0.10
 encapsulation dot1Q 10
 ip address 192.168.1.1 255.255.255.192
interface GigabitEthernet0/0.30
 encapsulation dot1Q 30
 ip address 192.168.1.65 255.255.255.240
\end{verbatim}

\section{OSPF y balanceo de costo igual}

OSPF es un protocolo de estado de enlace. Cada router forma adyacencias, construye una base de datos de estado de enlace y ejecuta SPF. Se utilizo un solo proceso OSPF 1 y area 0. Las interfaces hacia usuarios son pasivas: sus redes se anuncian, pero no se envian saludos OSPF a los hosts.

Para obligar dos caminos de costo igual entre los routers OSPF se asignaron los siguientes costos:
\begin{itemize}
\item Enlace directo R-OSPF-1--R-OSPF-2: costo 20.
\item Cada enlace de acceso a R-CENTRAL: costo 10.
\item Interfaz LAN de destino: costo 1.
\end{itemize}

Desde R-OSPF-1 hacia \texttt{192.168.0.0/24}:
\begin{calculo}
\[
C_{directo}=20+1=21.
\]
\[
C_{central}=10+10+1=21.
\]
Como $C_{directo}=C_{central}$, ambos siguientes saltos se instalan en la tabla: \texttt{10.0.0.2} y \texttt{10.0.0.6}. Esta es una comprobacion ECMP real, no solamente una afirmacion teorica.
\end{calculo}

\section{EIGRP, DUAL y balanceo desigual 75/25}

EIGRP usa DUAL. El sucesor es el mejor siguiente salto. Un vecino es sucesor factible cuando su distancia reportada (RD) es estrictamente menor que la distancia factible del sucesor (FD):
\[
RD_{alterno}<FD_{sucesor}.
\]
Con los valores $K_1=K_3=1$ y los demas $K=0$, la metrica clasica es
\[
M=256\left(\left\lfloor\frac{10^7}{BW_{min}}\right\rfloor+\sum\frac{delay_{\mu s}}{10}\right).
\]

Para R-EIGRP-1 hacia \texttt{172.16.0.0/25}, $BW_{min}=1544$ kbit/s:
\[
\left\lfloor\frac{10^7}{1544}\right\rfloor=6476.
\]

\subsection{Camino directo: sucesor}
El retardo total observado es $20010\ \mu s$, equivalente a 2001 unidades de decenas de microsegundo:
\begin{calculo}
\[
M_d=256(6476+2001)=256(8477)=\boxed{2\,170\,112}.
\]
\end{calculo}

\subsection{Camino por R-CENTRAL: candidato factible}
Se configuro \texttt{delay 18854} en la salida del router de acceso hacia el central y \texttt{delay 100} en la salida del central hacia el router remoto. El retardo total resulta $189550\ \mu s=18955$ unidades:
\begin{calculo}
\[
M_a=256(6476+18955)=256(25431)=\boxed{6\,510\,336}.
\]
La distancia reportada real del central fue
\[
RD_a=1\,683\,712<FD_d=2\,170\,112,
\]
por lo que se cumple la condicion de factibilidad y no existe riesgo de bucle segun DUAL.
\end{calculo}

\subsection{Variance y proporcion de trafico}
\[
\frac{M_a}{M_d}=\frac{6\,510\,336}{2\,170\,112}=3.
\]
Se necesita \texttt{variance 3}, porque $M_a\le3M_d$. Con reparto inversamente proporcional a la metrica:
\[
P_d=\frac{M_a}{M_d+M_a}=\frac{3}{1+3}=75\%,
\qquad
P_a=\frac{M_d}{M_d+M_a}=\frac{1}{4}=25\%.
\]
Los comandos relevantes son \texttt{variance 3}, \texttt{maximum-paths 2} y \texttt{traffic-share balanced}.

\section{Redistribucion y seleccion de OSPF E1}

R-CENTRAL ejecuta OSPF 1 y EIGRP 100; es el unico punto de redistribucion. Al llevar rutas OSPF hacia EIGRP es obligatoria una metrica semilla porque las rutas externas no poseen los cinco componentes de la metrica EIGRP:
\begin{verbatim}
redistribute ospf 1 metric 1544 20000 255 1 1500
\end{verbatim}
Los valores representan ancho de banda, retardo, confiabilidad, carga y MTU.

En sentido EIGRP hacia OSPF se eligio tipo E1. Debido a que el IOS emulado no acepta \texttt{summary-address} para rutas externas OSPF, se creo una ruta estatica de resumen hacia Null0 y se redistribuyo solamente esa ruta:
\begin{verbatim}
ip route 172.16.0.0 255.255.255.0 Null0
router ospf 1
 redistribute static metric-type 1 subnets
\end{verbatim}
R-CENTRAL conserva las rutas EIGRP especificas para reenviar trafico. Por coincidencia de prefijo mas largo, las rutas /25, /26, /27 y /28 ganan sobre la ruta /24 hacia Null0. OSPF recibe unicamente \texttt{172.16.0.0/24} como externa E1. Una ruta E1 suma el costo externo y el costo interno hasta el ASBR; una E2 conserva principalmente el costo externo. E1 refleja mejor la topologia y permite que el costo hacia R-CENTRAL influya en la seleccion.

El resumen OSPF hacia EIGRP se aplico en ambas interfaces EIGRP del central:
\begin{verbatim}
ip summary-address eigrp 100 192.168.0.0 255.255.254.0
\end{verbatim}
Los routers EIGRP reciben una sola ruta \texttt{192.168.0.0/23}.

\begin{decision}[title={Metodo de resumen compatible con Packet Tracer}]
La ruta Null0 no reemplaza la conectividad EIGRP: funciona como objeto de resumen estable en el ASBR. La tabla del central mantiene las rutas especificas aprendidas por EIGRP, mientras que OSPF publica una sola externa E1. Este metodo evita inyectar por separado las VLAN EIGRP y cumple el objetivo funcional de sumarizacion aun cuando el IOS 1941 emulado rechaza el comando OSPF \texttt{summary-address}.
\end{decision}

\section{Lista de comprobacion manual}
\begin{enumerate}
\item Revisar que cada interfaz tenga la IP y mascara de la tabla.
\item Verificar que los enlaces seriales esten \texttt{up/up} y que solo el lado DCE tenga reloj.
\item Confirmar dos vecinos OSPF en cada router del triangulo.
\item Confirmar dos vecinos EIGRP en cada router del triangulo.
\item Verificar dos siguientes saltos de costo 21 en OSPF.
\item Verificar las metricas 2,170,112 y 6,510,336 en EIGRP.
\item Comprobar $RD<FD$ y \texttt{variance 3}.
\item Confirmar la ruta resumida \texttt{192.168.0.0/23} en ambos routers EIGRP.
\item Confirmar una sola ruta OSPF E1 \texttt{172.16.0.0/24} en los routers OSPF.
\item Ejecutar pings y trazas entre todos los segmentos.
\end{enumerate}

\section{Fuentes de apoyo del curso}
\begin{itemize}
\item \emph{Lab3 - Algoritmos de Enrutamiento}: requisitos y rubrica del laboratorio.
\item \emph{Clase Direccionamiento IP - Subnetting UVG}: formula de hosts, CIDR y asignacion VLSM.
\item \emph{Enrutamiento Dinamico UVG}: OSPF, EIGRP, DUAL, metricas y distancias administrativas.
\item \emph{Guia de configuracion OSPF}: area 0, interfaces pasivas, enlaces seriales y comandos de verificacion.
\end{itemize}

\end{document}


